% !TeX root = ../../Book2.tex
\OptionalChapter{紧流形上的标量椭圆方程}{b2:ch:26}
\BookChapterNumber{b2:ch:26}

\begin{chapterabstract}
在紧致无边界的空间上，Poisson 方程的源项必须与常数方向相容，解也只能在固定平均值后唯一确定。本章从坐标图和正体积密度构造函数空间，用能量方法求解，再把欧氏内部估计拼接成全局正则性。移位逆算子的紧性给出光滑特征函数的正交展开；环面与球面展示这种展开如何化成可以逐项计算的公式。
\end{chapterabstract}

\begin{learninggoals}
\begin{itemize}
\item\label{b2:item:26-goal-1}理解体积密度的换元规律，利用有限坐标图定义 Sobolev 空间并解释其紧嵌入。
\item\label{b2:item:26-goal-2}从度量、梯度和散度推出 Laplace--Beltrami 算子，写出含变系数和非负势的弱方程。
\item\label{b2:item:26-goal-3}识别各连通分支上的零空间与相容条件，证明归一化解的存在、唯一性及稳定估计。
\item\label{b2:item:26-goal-4}用切断函数拼接内部正则性，辨认紧预解式、特征展开和自伴算子的定义域。
\item\label{b2:item:26-goal-5}在环面与球面上计算特征值及 Poisson 解，并比较共振、几何尺度和非紧空间中的变化。
\end{itemize}
\end{learninggoals}

设想热量在一个封闭曲面上传播，同时各处持续有热源或冷源。若希望达到稳态，总的输入和输出必须平衡；没有边界可供多余的热量流走。即使已经平衡，把温度整体升高一个常数也不改变热流，因而方程不能决定绝对温度。若空间由两个互不相通的部分组成，平衡和归一化还必须在每一部分单独考虑。

这些现象在方程 \(-\Delta_g u=f\) 中分别表现为源项的相容条件和解的常数自由度。下文先说明曲面上的面积、梯度与弱导数如何计算，再证明这些直观条件确实足以求解。正则性把能量解提升为光滑解，谱展开则进一步说明各种空间振荡各自怎样响应源项。全章的流形均紧致、光滑且无边界；定向和连通性在需要时另行说明。

关于局部坐标与单位分解，可参阅 Tu 的《An Introduction to Manifolds》\cite[Section 13]{Tu2011Introduction}。Taylor 的《Partial Differential Equations I: Basic Theory》\cite[Chapter 2, Section 4; Chapter 4, Section 3; Chapter 5, Section 1]{Taylor2011Partial}分别讨论 Riemann 流形上的 Laplace 算子、紧流形上的 Sobolev 空间以及椭圆问题的存在性与正则性；可按这些主题继续学习。

% !TeX root = ../../Book2.tex
\section{紧流形上的函数空间}
\label{b2:sec:2601}

在一块平面区域上，函数的可积性由 Lebesgue 测度衡量，弱导数则借助坐标方向定义。把区域换成球面后，函数仍然有大小和变化率，但一个坐标系只能描述球面的一部分。我们需要把各坐标片上的积分与导数组织起来，使换一组坐标之后得到同一个函数空间。紧性在这里有一个具体用途：有限个坐标片就足以覆盖整个空间，局部估计因而可以相加。

\subsection{坐标图与单位分解}
\label{b2:subsec:260101}

本章研究非空、紧致、无边界的光滑 Riemann 流形 \((M,g)\)，维数为整数 \(n\geq1\)。流形按第一卷附录 E 的约定是 Hausdorff、第二可数的；一张坐标图 \(x:U\to x(U)\subset\mathbb R^n\) 把开集 \(U\) 与欧氏开集光滑地识别，交叠处的坐标转换是光滑微分同胚。无边界表示所有点都有这样的欧氏开邻域，而不是半空间邻域。紧致无边界流形也常称为\term[bi-liuxing]{闭流形}[closed manifold]；这里的“闭”包含紧致这一要求。

度量 \(g\) 给每个切空间 \(T_pM\) 一个随 \(p\) 光滑变化的正定内积。在坐标基下，它由对称正定矩阵 \((g_{ij})\) 表示。第一卷定义 E.2.7 和定理 E.2.8 给出了这个概念及其由单位分解拼接的构造。球面采用欧氏内积在切空间上的限制，平坦环面采用从欧氏空间传下来的内积；同一个光滑流形也可以配上不同度量。

先记录一个与相容条件有关的事实。流形局部道路连通，所以每个连通分支都是开集：连通坐标小球中任意两点可以沿线段的逆像连接，故它包含在一个分支内。所有分支组成开覆盖，紧性使分支只有有限个，记为 \(M_1,\ldots,M_q\)。它们也都是闭集，因而各自紧致。函数 \(\mathbf1_{M_j}\) 局部恒定，特别是光滑函数。以后说“常数函数”时，需要辨认是整个 \(M\) 上同一个常数，还是每个分支各取一个常数。

\begin{lemma}[有限坐标局部化]\label{b2:lem:28-finite-localization}
设 \(M\) 为非空、紧致、无边界的 \(n\geq1\) 维光滑流形。存在有限个坐标图 \(x_a:U_a\to x_a(U_a)\)、开集 \(V_a\Subset U_a\) 及函数 \(\chi_a\in C^\infty(M;\mathbb R)\)，使
\[
 M=\bigcup_{a=1}^N V_a,\qquad
 0\leq\chi_a\leq1,\qquad
 \operatorname{supp}\chi_a\Subset U_a,\qquad
 \sum_{a=1}^N\chi_a=1.
\]
每个 \(U_a\) 还可以选成较大坐标图中的相对紧开集。对每个 \(a\)，存在 \(\eta_a\in C_c^\infty(U_a)\)，在 \(\operatorname{supp}\chi_a\) 的一个邻域上等于 \(1\)。
\end{lemma}
\begin{proof}
在任一点附近，先在一张图中选两个同心小球，其闭包都留在图内，把小球的逆像记为 \(V\)，较大球的逆像记为 \(U\)。所有这样的 \(V\) 覆盖 \(M\)，取有限子覆盖 \(V_a\)。第一卷定理 E.2.5 的光滑单位分解可用于这个有限坐标覆盖；也可以直接作如下构造。取 \(b_a\in C_c^\infty(U_a)\)，满足 \(0\leq b_a\leq1\)、\(b_a=1\) 于 \(\overline V_a\) 附近。在图外补零仍光滑，因为支集留在图内。总和 \(s=\sum_a b_a\) 处处至少为 \(1\)，故取 \(\chi_a=b_a/s\) 即可。最后，在图内对紧集 \(\operatorname{supp}\chi_a\) 再取一次光滑截断，得到 \(\eta_a\)。
\end{proof}

较大的坐标片为磨光和求导留出空间。直接把一个只在坐标域内定义的函数补零，通常会在域的边缘产生跳跃；先乘 \(\chi_a\)，则它在接近边缘时已经为零。图\ref{b2:fig:28-chart-localization}表示这个安排。

\begin{figure}[htbp]
\centering
\begin{tikzpicture}[scale=.8]
\draw[thick] (0,0) ellipse (3 and 1.65);
\draw[dashed] (0,0) ellipse (2.35 and 1.18);
\fill[blue!12] (0,0) ellipse (1.6 and .72);
\draw[blue!65!black] (0,0) ellipse (1.6 and .72);
\node at (0,0) {\(\operatorname{supp}\chi_a\)};
\node at (0,1.0) {\(\eta_a=1\)};
\node at (2.65,.6) {\(U_a\)};
\draw[->,thick] (3.35,0)--(5,0) node[midway,above] {\(x_a\)};
\draw[thick] (5.5,-1.65) rectangle (10.5,1.65);
\fill[blue!12] (8,0) ellipse (1.4 and .72);
\draw[blue!65!black] (8,0) ellipse (1.4 and .72);
\node at (8,0) {\((\chi_a u)\circ x_a^{-1}\)};
\node at (8,-1.2) {边缘附近恒为零};
\node at (8,1.95) {\(x_a(U_a)\subset\mathbb R^n\)};
\end{tikzpicture}
\caption{坐标局部化的支集安排。图形仅示意嵌套关系；右侧函数可补零到整个欧氏空间，再作磨光。}
\label{b2:fig:28-chart-localization}
\end{figure}

\subsection{积分与 Sobolev 空间}
\label{b2:subsec:260102}

度量首先给出积分所需的体积。若 \(x=(x^1,\ldots,x^n)\) 是局部坐标，记
\[
 \rho_x(x)=\sqrt{\det(g_{ij}(x))}>0.
\]
换成坐标 \(y\)，令 \(J=\partial x/\partial y\)，内积矩阵的变换公式为 \(g_y=J^{\mathsf T}g_xJ\)。取行列式及正平方根，得到
\begin{equation}\label{b2:eq:28-volume-change}
 \rho_y(y)=\rho_x(x(y))\,|\det J(y)|.
\end{equation}
这个绝对值与多重积分换元中的 Jacobi 因子完全相配。因此，局部表达式 \(\rho_x\,\mathrm dx\) 能够定义一个全局正测度。

\begin{definition}[Riemann 体积测度]\label{b2:def:28-volume}
设 \((M,g)\) 为非空、紧致、无边界的 \(n\geq1\) 维光滑 Riemann 流形。取\cref{b2:lem:28-finite-localization}中的有限坐标图和单位分解。对非负可测函数 \(f\)，规定
\begin{equation}\label{b2:eq:28-volume-integral}
 \int_M f\,\mathrm dV_g
 =\sum_{a=1}^N\int_{x_a(U_a)}
 (\chi_a f)\circ x_a^{-1}(x)\,\rho_{x_a}(x)\,\mathrm dx.
\end{equation}
由此得到的正测度称为 \(g\) 的\term[riemann-tiji-cedu]{Riemann 体积测度}[Riemannian volume measure]，记作 \(V_g\)。可测集 \(E\) 的体积也记为 \(\operatorname{Vol}_g(E)=V_g(E)\)。可测性按坐标中的 Lebesgue 可测性理解，零测集取相应完备化。
\end{definition}

\begin{proposition}[体积积分的选择无关性]\label{b2:prop:28-volume-invariant}
对非空、紧致、无边界的光滑 Riemann 流形 \((M,g)\)，定义\eqref{b2:eq:28-volume-integral}与有限坐标图、单位分解的选择无关。在任一坐标域 \(U\) 内，测度的表达式为 \(\mathrm dV_g=\rho_x\,\mathrm dx\)。它满足 \(0<V_g(M)<\infty\)，每个非空开集有正测度；在每个相对紧坐标子域上，它与 Lebesgue 测度互相具有有界的正密度。
\end{proposition}
\begin{proof}
若另选 \((W_b,y_b,\theta_b)\)，在原来的每一项中插入 \(\sum_b\theta_b=1\)。双重和的每项具有紧支集，留在 \(U_a\cap W_b\) 内；\eqref{b2:eq:28-volume-change}及欧氏换元公式使它等于用 \(y_b\) 坐标计算的同一项。再对 \(a\) 求和，\(\sum_a\chi_a=1\) 给出另一套定义。非负函数允许积分为无穷，有限求和和换元仍成立。

对集合的示性函数应用公式，欧氏测度的可数可加性及单调收敛给出一个测度。对支集位于某图中的函数，在上述比较中始终使用该图，便得到局部密度公式。光滑正函数 \(\rho_x\) 在坐标紧集上有正下界和有限上界；有限个这样的紧集覆盖 \(M\)，所以总体积有限。任一非空开集包含一个坐标小球，球上密度有正下界，故其体积正。
\end{proof}

对函数积分使用正体积密度，不需要选择定向。第一卷附录 E 中顶次微分形式的积分带有定向符号；这里的绝对值把体积当作非负大小。若 \(M\) 可定向并已选定向，\(\mathrm dV_g\) 也可写成与该定向相容的体积形式；改变定向不会改变当前的正测度。以后在局部作积分分部时，我们直接使用正密度公式。

对实值可测函数，按 \(V_g\) 定义 \(L^p(M)\)，\(1\leq p\leq\infty\)；函数仍按几乎处处相等识别。特别地，
\[
 (u,v)_{L^2(M)}=\int_M uv\,\mathrm dV_g.
\]
局部正密度保证坐标中的零测集与 \(V_g\)-零测集相同，光滑坐标转换也保持这类零测集。这使弱导数作用于函数等价类成为可能。

\begin{definition}[整数阶 Sobolev 空间]\label{b2:def:28-sobolev}
设 \((M,g)\) 为非空、紧致、无边界的 \(n\geq1\) 维光滑 Riemann 流形，\(k\in\mathbb N_0\)。固定\cref{b2:lem:28-finite-localization}中的有限局部化。对实值 \(u\in L^2(M)\)，把
\[
 u_a=(\chi_a u)\circ x_a^{-1}
\]
在 \(x_a(U_a)\) 外补零。若每个 \(u_a\) 属于 \(H^k(\mathbb R^n)=W^{k,2}(\mathbb R^n)\)，则称 \(u\in H^k(M)\)，并定义
\begin{equation}\label{b2:eq:28-chart-sobolev-norm}
 \|u\|_{H^k(M)}^2
 =\sum_{a=1}^N\sum_{|\alpha|\leq k}
 \|\partial^\alpha u_a\|_{L^2(\mathbb R^n)}^2.
\end{equation}
导数均为欧氏分布意义下的弱导数。复值版本对实部和虚部分别定义。
\end{definition}

这个范数包含零阶项。仅测量导数会漏掉每个连通分支上的常数，不能得到整个函数空间的范数。\(H^0(M)\) 与 \(L^2(M)\) 的范数等价，但通常不逐项相等：一方面密度在各支集上有界，另一方面
\[
 \frac1N\leq\sum_a\chi_a(p)^2\leq1
\]
来自 \(\sum_a\chi_a=1\) 和 Cauchy--Schwarz 不等式。

换坐标时要用一个具体估计。设 \(\Phi:O\to O'\) 是欧氏开集间的光滑微分同胚，\(b\in C_c^\infty(O)\)，\(h\in H^k_{\mathrm{loc}}(O')\)。在包含 \(\Phi(\operatorname{supp}b)\) 的一个紧邻域 \(K'\Subset O'\) 上，存在常数 \(C\) 使
\begin{equation}\label{b2:eq:28-coordinate-transfer}
 \|b\cdot(h\circ\Phi)\|_{H^k(O)}
 \leq C\sum_{|\beta|\leq k}\|\partial^\beta h\|_{L^2(K')}.
\end{equation}
这里可将左侧补零。对光滑 \(h\)，反复使用乘积和链式法则，左侧每个导数都是 \((\partial^\beta h)\circ\Phi\) 与光滑系数的有限和，且 \(|\beta|\leq k\)。系数在固定紧集上有界；换元公式中的 \(|\det D\Phi^{-1}|\) 也有界，因此每项满足所写估计。对一般 \(h\)，先在 \(K'\) 附近乘一个等于 \(1\) 的截断，再磨光。第一卷命题22.1.2保证全部不超过 \(k\) 阶的弱导数在 \(L^2\) 中收敛；将光滑估计用于差，再利用弱导数的闭性即可取极限。这也是第一卷定理21.3.7之后高阶坐标变换论证在紧支集上的形式。

\begin{theorem}[Sobolev 空间的基本性质]\label{b2:thm:28-sobolev-basic}
设 \((M,g)\) 为非空、紧致、无边界的 \(n\geq1\) 维光滑 Riemann 流形，\(k\in\mathbb N_0\)，函数取实值或复值。空间 \(H^k(M)\) 不依赖有限坐标局部化的选择，不同选择给出等价范数。它是 Hilbert 空间，\(C^\infty(M)\) 在其中稠密。一个 \(u\in L^2(M)\) 属于 \(H^k(M)\)，当且仅当它在每张坐标图内局部属于 \(H^k\)。
\end{theorem}
\begin{proof}
先说明局部性质。若所有 \(u_a\in H^k\)，则在任意坐标紧子域内用 \(u=\sum_a\chi_a u\) 分解。每一项通过坐标转换表达；在交叠部分取有限个较小坐标邻域和截断，\eqref{b2:eq:28-coordinate-transfer}保证其 \(H^k\) 性质。所需邻域只有有限个，因为被考察的支集紧含于交叠区域。反之，若 \(u\) 在各图内局部属于 \(H^k\)，则乘 \(\chi_a\) 后的支集留在内部，乘积法则和补零给出 \(u_a\in H^k(\mathbb R^n)\)。

换成另一套 \((W_b,y_b,\theta_b)\)，将每个 \(\theta_bu\) 写成 \(\sum_a\theta_b\chi_au\)。每个交叠项的支集包含于 \(\operatorname{supp}\theta_b\cap\operatorname{supp}\chi_a\Subset W_b\cap U_a\)。取在该支集附近等于 \(1\) 的截断并用\eqref{b2:eq:28-coordinate-transfer}，得到新范数不超过原范数的常数倍。交换两套局部化就得到反向估计。

范数\eqref{b2:eq:28-chart-sobolev-norm}来自各弱导数的 \(L^2\) 内积之和。设 \((u_j)\) 是此范数下的 Cauchy 列。零阶控制使它在 \(L^2(M)\) 中趋于某个 \(u\)，而对每个 \(a\)，\((\chi_au_j)\circ x_a^{-1}\) 在 \(H^k(\mathbb R^n)\) 中趋于某个 \(v_a\)。在 \(L^2\) 中辨认极限，得到 \(v_a=(\chi_au)\circ x_a^{-1}\)，在图外为零。因此 \(u\in H^k(M)\)，所有有限个坐标分量同时收敛，故 \(u_j\to u\) 于 \(H^k(M)\)。

最后证明稠密性。每个 \(u_a\) 有固定内部紧支集。用紧支撑磨光核取 \(u_{a,\varepsilon}=\varrho_\varepsilon*u_a\)；当 \(\varepsilon\) 足够小时，所有支集仍留在同一个紧子集 \(K_a\Subset x_a(U_a)\)，且 \(u_{a,\varepsilon}\to u_a\) 于 \(H^k(\mathbb R^n)\)。将各 \(u_{a,\varepsilon}\circ x_a\) 补零并相加，得到光滑全局函数 \(v_\varepsilon\)。对其与 \(u\) 的差再乘任意 \(\chi_b\)，使用有限个坐标转换估计，便有
\[
 \|v_\varepsilon-u\|_{H^k(M)}
 \leq C\sum_a\|u_{a,\varepsilon}-u_a\|_{H^k(\mathbb R^n)}\longrightarrow0.
\]
这同时说明各局部磨光结果在交叠处不必相等：相加的是全局补零后的光滑函数，极限由 \(\sum_a\chi_a u=u\) 辨认。
\end{proof}

例如，圆 \(\mathbb R/(2\pi\mathbb Z)\) 上的函数可看作 \(2\pi\)-周期函数。通过跨过角度 \(0\) 的另一张图，\(H^1\) 的定义要求两端相接；仅在区间 \((0,2\pi)\) 内有平方可积导数还不够。角度函数 \(u(\theta)=\theta\) 在这个区间内光滑，但作为圆上的函数在接缝处跳跃，因而不属于圆上的 \(H^1\)。局部坐标要覆盖接缝，正是空间定义能够识别这类现象的原因。

\subsection{稠密性与紧嵌入}
\label{b2:subsec:260103}

有了局部化，欧氏空间中的紧性可以搬到 \(M\) 上。关键是各局部函数拥有共同的有界支集：有界列不能把质量无限移到远处。

\begin{theorem}[紧流形上的 Rellich 紧嵌入]\label{b2:thm:28-manifold-rellich}
设 \((M,g)\) 为非空、紧致、无边界的 \(n\geq1\) 维光滑 Riemann 流形，\(k\in\mathbb N_0\)，函数取实值或复值。嵌入
\[
 H^{k+1}(M)\hookrightarrow H^k(M)
\]
连续且紧。特别地，每个 \(H^1(M)\) 有界列都有在 \(L^2(M)\) 中强收敛的子列。
\end{theorem}
\begin{proof}
连续性由定义中的导数项直接得到。设 \((u_j)\) 在 \(H^{k+1}(M)\) 中有界。对每个 \(a\) 和 \(|\alpha|\leq k\)，局部函数
\[
 \partial^\alpha\bigl((\chi_a u_j)\circ x_a^{-1}\bigr)
\]
在 \(H^1(\mathbb R^n)\) 中有界，并且都支撑在同一个坐标紧集。选一个包含该紧集的大球，它们的限制属于 \(H_0^1\)；这可由内部磨光逼近直接看出。第一卷推论25.2.6给出该有界球上 \(H_0^1\hookrightarrow L^2\) 的紧性。依次抽取子列，使有限个 \((a,\alpha)\) 的分量全部在 \(L^2\) 中收敛。于是所选子列在\eqref{b2:eq:28-chart-sobolev-norm}中为 Cauchy 列，\cref{b2:thm:28-sobolev-basic}的完备性给出所求极限。
\end{proof}

正则性提高到一定阶数后，函数还会具有经典导数。取有限图册中的紧子域，定义 \(C^r\) 范数为这些子域上不超过 \(r\) 阶坐标导数的上确界之和；不同图册的范数等价。

\begin{proposition}[连续可微代表]\label{b2:prop:28-sobolev-continuous}
设 \((M,g)\) 为非空、紧致、无边界的 \(n\geq1\) 维光滑 Riemann 流形，\(k,r\in\mathbb N_0\) 且 \(k>r+n/2\)。每个实值或复值 \(u\in H^k(M)\) 有唯一的 \(C^r(M)\) 代表，并有
\[
 \|u\|_{C^r(M)}\leq C\|u\|_{H^k(M)}.
\]
因而一个属于所有整数阶 \(H^k(M)\) 的函数等价类具有唯一光滑代表。
\end{proposition}
\begin{proof}
每个局部补零函数有固定有限测度支集。可选 \(1<p<2\) 足够接近 \(2\)，使 \(k-n/p>r\) 且 \(k-n/p\) 不是整数。Hölder 不等式给出局部 \(W^{k,p}\) 范数受 \(H^k\) 范数控制。第一卷定理24.4.3对全空间函数给出 \(W^{k,p}\) 的连续可微代表，其可微阶数为 \(\lfloor k-n/p\rfloor\geq r\)，并控制对应范数。这里通过稍减 \(p\) 保留了严格的指数余量，因此也覆盖 \(k-n/2\) 为整数的情形。

将这些局部代表拉回并相加，得到 \(u\) 的全局 \(C^r\) 代表和所写估计。若两个连续代表几乎处处相同而在一点不同，连续性将使它们在一个非空开集上不同，违反该开集有正体积。故代表唯一。让 \(r\) 任意增大，所有代表因唯一性相同，得到光滑性。
\end{proof}

在弱解的唯一性证明中，还会出现“导数为零意味着常数”这个结论。对弱导数需要给出与经典微积分衔接的步骤。

\begin{lemma}[弱导数为零的函数]\label{b2:lem:28-zero-gradient}
设 \(M\) 为非空、紧致、无边界的光滑流形，配有光滑 Riemann 度量，\(u\in H^1(M;\mathbb R)\)。若 \(u\) 在各坐标图中的全部一阶弱导数为零，则它在每个连通分支上几乎处处等于一个常数。反之，分支常数函数属于 \(C^\infty(M)\)，且所有这些弱导数均为零。
\end{lemma}
\begin{proof}
在任一坐标球的紧内部小球上磨光 \(u\)。弱导数与磨光交换，故磨光函数的梯度恒为零。沿小球内线段积分，它是常数。磨光函数在 \(L^2\) 中趋于 \(u\)，而常数函数构成 \(L^2\) 中的闭子空间，所以 \(u\) 在小球上几乎处处为常数。两个这样的球若相交，交叠开集有正体积，两个常数必须相等。每个连通分支道路连通，一条连接路径的紧像可以用有限条交叠小球链覆盖，故常数在整个分支一致。分支数量有限且各分支开闭，反向结论随即成立。
\end{proof}

以后用 \(H^{-1}(M)\) 表示实 Hilbert 空间 \(H^1(M)\) 的连续对偶，配对记作 \(\langle F,v\rangle\)。任何 \(f\in L^2(M)\) 通过 \(\langle f,v\rangle=\int_Mfv\,\mathrm dV_g\) 给出其中一个元素。这样的数据可以直接作用在弱解的测试函数上；它比要求右端逐点有定义更适合能量方法。下一节将用度量把一阶弱导数写成全局梯度，并将当前坐标范数与能量积分联系起来。

% !TeX root = ../../Book2.tex
\section{Laplace–Beltrami 算子与弱形式}
\label{b2:sec:2602}

待撰写。

\subsection{度量、梯度与散度}
\label{b2:subsec:260201}

待撰写。

\subsection{Laplace–Beltrami 算子的局部表达}
\label{b2:subsec:260202}

待撰写。

\subsection{能量型与弱方程}
\label{b2:subsec:260203}

待撰写。

% !TeX root = ../../Book2.tex
\section{标量椭圆问题的存在性与唯一性}
\label{b2:sec:2603}

待撰写。

\subsection{强制性与 Lax–Milgram 方法}
\label{b2:subsec:260301}

待撰写。

\subsection{零空间与相容条件}
\label{b2:subsec:260302}

待撰写。

\subsection{平均值归一化与唯一性}
\label{b2:subsec:260303}

待撰写。

% !TeX root = ../../Book2.tex
\section{局部估计与全局正则性}
\label{b2:sec:2604}

待撰写。

\subsection{坐标局部化与估计拼接}
\label{b2:subsec:260401}

待撰写。

\subsection{二阶 Sobolev 估计}
\label{b2:subsec:260402}

待撰写。

\subsection{高阶正则性与光滑性}
\label{b2:subsec:260403}

待撰写。

% !TeX root = ../../Book2.tex
\section{紧预解式与离散谱}
\label{b2:sec:2605}

待撰写。

\subsection{移位逆算子的构造与紧性}
\label{b2:subsec:260501}

待撰写。

\subsection{特征值与正交完备性}
\label{b2:subsec:260502}

待撰写。

\subsection{谱展开与椭圆方程}
\label{b2:subsec:260503}

待撰写。

% !TeX root = ../../Book2.tex
\section{典型模型与理论边界}
\label{b2:sec:2606}

特征函数展开把紧流形上的方程化成一列标量方程。若流形具有较多对称性，还能把这些特征函数具体写出来：平坦环面的平移对称性产生 Fourier 指数，球面的旋转对称性产生球谐函数。比较两个模型，可以看见频率、重数、体积归一化与几何尺度怎样进入同一个解公式。

\subsection{平坦环面上的 Fourier 模型}
\label{b2:subsec:260601}

固定整数 \(n\geq1\)，取
\[
 \mathbb T_{2\pi}^n\coloneqq\mathbb R^n/(2\pi\mathbb Z)^n.
\]
它把每个坐标相差 \(2\pi\) 的点识别起来；其光滑函数等同于在每个变量上都以 \(2\pi\) 为周期的光滑函数。欧氏度量在平移下不变，因而给出商空间上的平坦度量。相应的体积积分是在基本区域 \([0,2\pi]^n\) 上的 Lebesgue 积分，体积为 \((2\pi)^n\)，并且
\[
 \Delta_g=\sum_{a=1}^n\partial_{x_a}^2.
\]
基本区域的相对面已经识别，环面没有边界；对周期函数分部积分时，相对面的边界项相消。

\begin{proposition}[平坦环面的 Fourier 谱]\label{b2:prop:26-torus-spectrum}
设整数 \(n\geq1\)，\(\mathbb T_{2\pi}^n\) 取欧氏平坦度量和未归一化体积测度 \(\mathrm dx\)。则
\[
 E_k(x)=(2\pi)^{-n/2}e^{\mathrm i k\cdot x}\qquad(k\in\mathbb Z^n)
\]
组成复 \(L^2(\mathbb T_{2\pi}^n,\mathrm dx)\) 的正交归一基，满足
\[
 -\Delta_gE_k=|k|^2E_k.
\]
若对 \(f\in L^2(\mathbb T_{2\pi}^n;\mathbb C)\) 定义
\[
 \widehat f(k)=(2\pi)^{-n}\int_{[0,2\pi]^n}f(x)e^{-\mathrm ik\cdot x}\,\mathrm dx,
\]
则
\[
 f=\sum_{k\in\mathbb Z^n}\widehat f(k)e^{\mathrm ik\cdot x}\quad\text{于 }L^2,
 \qquad
 \|f\|_2^2=(2\pi)^n\sum_{k\in\mathbb Z^n}|\widehat f(k)|^2.
\]
特征值 \(\lambda\) 的重数是满足 \(|k|^2=\lambda\) 的整数向量个数，零特征空间为常数空间。
\end{proposition}
\begin{proof}
逐变量积分给出
\[
 \int_{[0,2\pi]^n}E_k\overline{E_\ell}\,\mathrm dx
   =\prod_{a=1}^n\left(\frac1{2\pi}\int_0^{2\pi}e^{\mathrm i(k_a-\ell_a)x_a}\,\mathrm dx_a\right)
   =\mathbf1_{\{k=\ell\}}.
\]
特征方程直接由二阶求导得到。

完备性的一维情形由\cref{b2:thm:02-l2-convergence}经换元 \(x=2\pi t\) 得到；其中原先的圆周为 \(\mathbb R/\mathbb Z\)，测度总质量为一，这里归一化因子正是换元后的调整。现按维数归纳。设 \(n\geq2\)，且 \(f\in L^2([0,2\pi]^n)\) 与全部 \(E_k\) 正交。将变量写成 \((x',t)\)。对每个 \(m\in\mathbb Z\)，令
\[
 g_m(x')=\int_0^{2\pi}f(x',t)e^{-\mathrm imt}\,\mathrm dt.
\]
Fubini 与 Cauchy--Schwarz 给出 \(g_m\in L^2([0,2\pi]^{n-1})\)。它的全部 \((n-1)\) 维 Fourier 系数为零，故归纳假设说明 \(g_m=0\) 几乎处处。整数 \(m\) 只有可数多个，可同时除去这些例外集。对剩下的几乎每个 \(x'\)，切片 \(f(x',\cdot)\in L^2(0,2\pi)\) 的所有一维系数为零，一维完备性给出该切片为零。因此 \(f=0\) 几乎处处，完备性成立。

正交基的系数为 \((f,E_k)=(2\pi)^{n/2}\widehat f(k)\)，代入 Parseval 就得到展开及范数等式。若 \(-\Delta_gf=\lambda f\)、\(f\in H^2\)，逐个取系数得到 \((|k|^2-\lambda)\widehat f(k)=0\)，所以整个特征空间恰由满足所列条件的指数函数张成。这样的整数向量处在有界球内，只有有限多个。
\end{proof}

实值函数满足 \(\widehat f(-k)=\overline{\widehat f(k)}\)，把 \(k\) 与 \(-k\) 配成一对，就得到实正弦、余弦基。正交归一基与通常 Fourier 系数之间的常数差值得保留：使用 \(E_k\) 时系数直接满足无额外因子的 Parseval，使用 \(e^{\mathrm ik\cdot x}\) 时则出现体积因子 \((2\pi)^n\)。

\ProseParagraphBreak

对平均值为零的 \(f\in L^2(\mathbb T_{2\pi}^n)\)，唯一平均值为零的 Poisson 解是
\[
 u=\sum_{k\ne0}\frac{\widehat f(k)}{|k|^2}e^{\mathrm ik\cdot x},
\]
级数在 \(H^2\) 中收敛。零频对应整体平均值，所以要单独去掉；其余每个频率都被除以长度的平方。例如在二维环面上，取
\[
 f(x_1,x_2)=\cos(2x_1-x_2)+2\sin x_2.
\]
两个频率分别为 \((2,-1)\) 与 \((0,1)\)，故解为
\[
 u(x_1,x_2)=\frac15\cos(2x_1-x_2)+2\sin x_2.
\]
第一项沿 \((2,-1)\) 方向振荡，频率平方为 \(5\)；第二项的频率平方为 \(1\)。逐项求导可以核验 \(-\Delta_g u=f\)。再用三角函数正交性，得到
\[
 \int_{\mathbb T_{2\pi}^2}|\nabla_g u|^2\,\mathrm dx
 =\int_{\mathbb T_{2\pi}^2}fu\,\mathrm dx
 =(2\pi)^2\left(\frac1{10}+2\right).
\]
这个计算同时检查了频率分母、体积因子与能量符号。

\subsection{球面上的球谐模型}
\label{b2:subsec:260602}

单位球面 \(S^n\subset\mathbb R^{n+1}\) 取欧氏空间诱导的度量，体积测度就是通常面积测度 \(\mathrm d\sigma\)。球面上的频率用多项式的齐次次数来组织。令 \(\mathcal H_\ell\) 为 \(\mathbb R^{n+1}\) 上 \(\ell\) 次齐次调和复多项式空间，即
\[
 P(tx)=t^\ell P(x),\qquad \Delta_{\mathbb R^{n+1}}P=0.
\]
将这些多项式限制在 \(S^n\) 上，得到 \(\ell\) 阶球谐空间 \(\mathcal Y_\ell\)。齐次性使限制映射单射：若 \(P\) 在单位球面上为零，则 \(P(r\omega)=r^\ell P(\omega)=0\)，于是它在整个空间为零。

\begin{proposition}[单位球面的球谐谱]\label{b2:prop:26-sphere-spectrum}
设整数 \(n\geq1\)，\(S^n\) 取单位球面的标准度量和面积测度。在每个 \(\ell\geq0\) 的球谐空间 \(\mathcal Y_\ell\) 中选取正交归一基 \((Y_{\ell,a})_{1\leq a\leq N_\ell}\)。则这些基向量合起来是复 \(L^2(S^n,\mathrm d\sigma)\) 的正交归一基，并满足
\[
 -\Delta_{S^n}Y_{\ell,a}=\ell(\ell+n-1)Y_{\ell,a},
\]
其中
\begin{equation}\label{b2:eq:26-spherical-multiplicity}
 N_\ell=\binom{\ell+n}{n}-\binom{\ell+n-2}{n}.
\end{equation}
当 \(\ell=0,1\) 时第二项按零处理。这些数 \(\ell(\ell+n-1)\) 是全部特征值，重数为 \(N_\ell\)。特别地，\(N_0=1\)、\(N_1=n+1\)，最小正特征值为 \(n\)。
\end{proposition}
\begin{proof}
在\cref{b2:prop:B2-polar-laplacian}中取环境维数 \(d=n+1\)，得到
\[
 \Delta_{\mathbb R^{n+1}}=\partial_r^2+\frac nr\partial_r+\frac1{r^2}\Delta_{S^n}.
\]
其中切向梯度来自诱导度量，角向散度与本章定义一致。代入 \(P(r\omega)=r^\ell Y(\omega)\)，得到
\[
 0=\Delta P=r^{\ell-2}\bigl(\ell(\ell+n-1)Y+\Delta_{S^n}Y\bigr),
\]
因而有特征关系。\cref{b2:prop:B2-harmonic-decomposition}给出齐次多项式的直和分解
\[
 \mathcal P_\ell=\mathcal H_\ell\oplus|x|^2\mathcal P_{\ell-2}.
\]
空间 \(\mathcal P_\ell\) 的单项式基由 \(n+1\) 个非负整数指数之和等于 \(\ell\) 来计数，维数为 \(\binom{\ell+n}{n}\)；维数相减便得到\eqref{b2:eq:26-spherical-multiplicity}。

相邻特征值之差为 \(2\ell+n>0\)，不同阶球谐因 Green 公式而正交。\cref{b2:thm:B2-spherical-expansion}已在面积测度下证明这些球谐的完备性：调和分解把球面多项式变成有限球谐和，多项式一致逼近连续函数，再由连续函数的 \(L^2\) 稠密性得到全部平方可积函数的展开。若 \(u\in H^2(S^n)\) 是任意特征函数，把方程与每个 \(Y_{\ell,a}\) 配对，就有
\[
 \bigl(\ell(\ell+n-1)-\lambda\bigr)(u,Y_{\ell,a})_{L^2}=0.
\]
完备性说明 \(\lambda\) 必等于某一个所列特征值，且 \(u\) 恰在对应的球谐空间中。因此列出的重数也就是整个特征空间的维数。
\end{proof}

当 \(n=1\) 时，球面就是单位圆，\(\lambda_\ell=\ell^2\)，且 \(N_0=1\)、\(N_\ell=2\)（\(\ell\geq1\)），对应同一频率下的正弦、余弦。当 \(n=2\) 时，\(\lambda_\ell=\ell(\ell+1)\)、\(N_\ell=2\ell+1\)。同一个阶数通常含多个独立方向，不能只用轴对称的 Legendre 多项式展开全部球面函数。

\ProseParagraphBreak

低阶源项可以直接计算。记 \(\omega=(\omega_1,\ldots,\omega_{n+1})\in S^n\)。一次函数 \(x_1\) 调和，故 \(\omega_1\) 的特征值为 \(n\)。二次多项式
\[
 x_1^2-\frac{|x|^2}{n+1}
\]
的欧氏 Laplace 为 \(2-2(n+1)/(n+1)=0\)，故其球面限制 \(\omega_1^2-1/(n+1)\) 的特征值为 \(2(n+1)\)。两者的平均值都为零：第一项由反射对称性得到，第二项由各坐标平方的积分相等及 \(\sum_a\omega_a^2=1\) 得到。因此
\[
 f(\omega)=\omega_1+3\left(\omega_1^2-\frac1{n+1}\right)
\]
的唯一零均值 Poisson 解为
\[
 u(\omega)=\frac{\omega_1}{n}
       +\frac{3}{2(n+1)}\left(\omega_1^2-\frac1{n+1}\right).
\]
一般平方可积零均值数据的解则是
\[
 u=\sum_{\ell\geq1}\sum_{a=1}^{N_\ell}
       \frac{(f,Y_{\ell,a})_{L^2}}{\ell(\ell+n-1)}Y_{\ell,a},
\]
由\cref{b2:thm:26-spectral-solvability}在 \(H^2(S^n)\) 中收敛。

\ProseParagraphBreak

几何尺度也能直接读出。对任意紧流形 \((M,g)\) 及常数 \(R>0\)，把度量换成 \(g_R=R^2g\)。坐标矩阵的逆乘以 \(R^{-2}\)，体积密度乘以 \(R^n\)，故
\[
 \mathrm dV_{g_R}=R^n\,\mathrm dV_g,\qquad
 \Delta_{g_R}=R^{-2}\Delta_g.
\]
若 \(e_j\) 对原度量正交归一，则对新度量相应的归一化特征函数为 \(R^{-n/2}e_j\)，特征值为 \(R^{-2}\lambda_j\)。半径为 \(R\) 的圆球面因此有特征值 \(R^{-2}\ell(\ell+n-1)\)。空间尺度越大，同一种形状的振荡变化越慢，相应的梯度能量也越小。

\subsection{标量理论与进一步推广的边界}
\label{b2:subsec:260603}

带边界的紧流形还需要规定边界上的行为。以区间 \([0,\pi]\) 为例，同一个微分表达式 \(-u''\)，配上 \(u(0)=u(\pi)=0\) 后，特征函数为 \(\sin(kx)\)、特征值 \(k^2\)（\(k\geq1\)）；配上 \(u'(0)=u'(\pi)=0\) 后，出现常数模态及 \(\cos(kx)\)、特征值 \(k^2\)（\(k\geq0\)）。这两个算子的定义域分别为
\[
 \{\,u\in H^2(0,\pi)\mid u(0)=u(\pi)=0\,\},\qquad
 \{\,u\in H^2(0,\pi)\mid u'(0)=u'(\pi)=0\,\}.
\]
端点条件使分部积分的边界项消失，却产生不同的零空间与 Poisson 相容条件。一般带边界流形上，除边界条件以外，还要检查边界的光滑性和相应的边界正则性估计，才能确定算子域与解的导数。

\ProseParagraphBreak

非紧流形则可能失去紧预解式。最直接的模型是 \(\mathbb R^n\)，\(n\geq1\)，及复 \(L^2(\mathbb R^n)\) 上定义域为 \(H^2(\mathbb R^n)\) 的算子 \(-\Delta\)。按本卷 Fourier 变换的 \(e^{-2\pi\mathrm ix\cdot\xi}\) 约定，\cref{b2:thm:04-plancherel-extension}与\cref{b2:thm:04-weak-derivative}把它变成乘法算子
\[
 \widehat{-\Delta u}(\xi)=4\pi^2|\xi|^2\widehat u(\xi).
\]
若 \(-\Delta u=\lambda u\) 且 \(u\in H^2\)，则 \(\widehat u\) 只能支撑在 \(4\pi^2|\xi|^2=\lambda\) 上。这是零测集（或空集），故平方可积的 \(\widehat u\) 只能为零。因此这个算子没有非零 \(L^2\) 特征函数。

\ProseParagraphBreak

尽管如此，每个 \(\lambda\geq0\) 都是谱点。取满足 \(4\pi^2|\xi_0|^2=\lambda\) 的 \(\xi_0\)，再取支撑在 \(\xi_0\) 的足够小邻域内、\(L^2\) 范数为一的光滑函数 \(h_\varepsilon\)，使其支集上
\(\bigl|4\pi^2|\xi|^2-\lambda\bigr|<\varepsilon\)。令 \(u_\varepsilon=\mathcal F^{-1}h_\varepsilon\)，则 \(u_\varepsilon\in H^2\)，且
\[
 \|u_\varepsilon\|_2=1,\qquad
 \|(-\Delta-\lambda)u_\varepsilon\|_2\leq\varepsilon.
\]
若存在有界逆，就会得到 \(1\leq C\varepsilon\)，矛盾。另一方面，对 \(z\notin[0,\infty)\)，乘子 \((4\pi^2|\xi|^2-z)^{-1}\) 定义从 \(L^2\) 到 \(H^2\) 的有界逆。故这里的谱是整个 \([0,\infty)\)。

\ProseParagraphBreak

这些谱点具有\term[lianxu-pu]{连续谱}[continuous spectrum]的特征：算子 \(-\Delta-\lambda\) 单射，值域稠密，却不是满射。单射性已由零测集论证得到。对任意频率数据，截去 \(|\xi|>N\) 和 \(\bigl|4\pi^2|\xi|^2-\lambda\bigr|<1/N\) 的区域，所余数据可以除以乘子得到 \(H^2\) 原像，而且这些截断在 \(L^2\) 中趋于原数据，所以值域稠密。非满射性也可直接构造：令
\[
 B_k=\{\,\xi\in\mathbb R^n\mid 2^{-k-1}<4\pi^2|\xi|^2-\lambda<2^{-k}\,\},
 \qquad h_k=|B_k|^{-1/2}\cdot\mathbf1_{B_k}\quad(k\geq1).
\]
各 \(B_k\) 有限正测度且两两不交，故 \(h=\sum_{k\geq1}2^{-k}h_k\) 属于 \(L^2\)。但 \(h/(4\pi^2|\xi|^2-\lambda)\) 在每个 \(B_k\) 上的平方积分至少为一，因此不属于 \(L^2\)。源项 \(\mathcal F^{-1}h\) 没有 \(H^2\) 解，这就核实了非满射性。

\ProseParagraphBreak

紧性失效也可在位置空间看见。取一个非零紧支集光滑函数 \(\phi\)，沿彼此越来越远的位置作平移 \(\phi_j\)，使支集两两不交。它们的 \(H^1\) 范数相同，而 \(\|\phi_j-\phi_k\|_2^2=2\|\phi\|_2^2\)（\(j\ne k\)），所以没有 \(L^2\) 收敛子列。紧流形上的 Rellich 结论因这种向无穷远平移而失效。更直接地，令 \(f_j=(-\Delta+1)\phi_j\)，则 \((f_j)\) 在 \(L^2\) 中有界，但 \((-\Delta+1)^{-1}f_j=\phi_j\) 没有收敛子列；这个预解式确实不紧。

\ProseParagraphBreak

标量函数还可以推广为向量场或微分形式。对 \(k\) 次微分形式，外微分 \(\mathrm d\) 将次数升一。在紧致无边界 Riemann 流形上，以度量和体积密度定义形式的 \(L^2\) 内积，局部分部积分得到将次数降一的形式伴随 \(\delta\)，其含义是 \((\mathrm d\alpha,\beta)_{L^2}=(\alpha,\delta\beta)_{L^2}\)。相应的\term[hodge-laplace-suanzi]{Hodge Laplace 算子}[Hodge Laplacian]为
\[
 \Delta_{\mathrm H}=\mathrm d\delta+\delta\mathrm d.
\]
在零次形式上，它与本章的非负算子 \(-\Delta_g\) 一致；在高次形式上，各个分量随坐标变换相互联系，零空间也要由 \(\mathrm d\omega=0\)、\(\delta\omega=0\) 两个条件共同描述。由此继续研究，需要建立形式上的 Sobolev 空间、椭圆系统估计以及闭形式的全局结构。标量情形中已经看到的分支常数、局部正则性与紧逆算子，是这些问题最具体的起点。


\begin{chaptersummary}
正体积密度使函数积分可以跨过坐标交叠；光滑单位分解使弱导数、逼近和紧性能够由有限个欧氏问题拼接。度量随后将导数写成梯度，把 Poisson 方程转化为能量恒等式。

无边界紧流形上的零梯度函数在每个连通分支上为常数。相容条件恰好消去源项在这些方向上的分量，平均值归一化则消去解在同一方向上的自由度。对非负势，只有势恒为零的分支仍保留常数零空间。

局部二阶估计与切断交换公式共同给出全局导数控制。它们既说明弱解的正则性，也保证特征向量具有光滑代表。紧解算子产生离散谱；一阶能量对应系数权重 \(1+\lambda_j\)，二阶算子域对应 \(1+\lambda_j^2\)。环面的整数频率与球面的球谐次数是这套结构的两个显式模型。边界条件、非紧性和未知量类型的变化，则分别改变定义域、紧性及零空间的含义。
\end{chaptersummary}

\BookExercises

% 本章八题均为自拟；题目与解答在下文依训练内容排列。

% 自拟：度量的共形改变与二维能量，训练坐标密度和逆度量的共同作用。
\Needspace{6\baselineskip}
\begin{exercise}
\label{b2:ex:26-conformal-change}
设 \((M,g)\) 为非空、紧致、无边界的 \(n\geq1\) 维光滑 Riemann 流形，\(\omega\in C^\infty(M;\mathbb R)\)，令新度量 \(\widetilde g=e^{2\omega}g\)。证明对 \(u\in C^2(M;\mathbb R)\)，有
\[
 \mathrm dV_{\widetilde g}=e^{n\omega}\mathrm dV_g,
 \quad \Delta_{\widetilde g}u
 =e^{-2\omega}\bigl(\Delta_g u
     +(n-2)g(\nabla_g\omega,\nabla_g u)\bigr).
\]
再证明当 \(n=2\) 时，\(u\in H^1(M)\) 的梯度能量对这样的度量改变保持不变。若 \(n=2\)、\(f\in C^\infty(M;\mathbb R)\)，把 \(-\Delta_{\widetilde g}u=f\) 改写为相对于 \(g\) 的弱方程，并写出每个连通分支上的相容条件。
\end{exercise}
\begin{solution}
局部矩阵满足 \(\widetilde g_{ij}=e^{2\omega}g_{ij}\)、\(\widetilde g^{ij}=e^{-2\omega}g^{ij}\)，故
\(\widetilde\rho=e^{n\omega}\rho\)。代入 Laplace--Beltrami 的坐标公式，得到
\[
 \begin{aligned}
 \Delta_{\widetilde g}u
 &=e^{-n\omega}\rho^{-1}
     \partial_i\bigl(e^{(n-2)\omega}\rho g^{ij}\partial_ju\bigr)\\
 &=e^{-2\omega}\Delta_gu
   +(n-2)e^{-2\omega}g^{ij}(\partial_i\omega)(\partial_ju),
 \end{aligned}
\]
最后一项就是题中两个梯度的内积。梯度长度平方按 \(e^{-2\omega}\) 缩放，所以
\[
 \int_M|\nabla_{\widetilde g}u|_{\widetilde g}^2\,\mathrm dV_{\widetilde g}
 =\int_Me^{(n-2)\omega}|\nabla_gu|_g^2\,\mathrm dV_g.
\]
先对光滑函数得到等式；两种光滑度量在紧流形上的 \(H^1\) 范数等价，光滑逼近遂将它推广到 \(H^1\)。当 \(n=2\) 时权重等于 \(1\)，就是能量不变性。

二维时，原方程等价于 \(-\Delta_gu=e^{2\omega}f\)，其弱式为
\[
 \int_Mg(\nabla_gu,\nabla_gv)\,\mathrm dV_g
 =\int_Me^{2\omega}fv\,\mathrm dV_g
 \quad(v\in H^1(M)).
\]
因此每个连通分支 \(M_j\) 上的相容条件为
\(\int_{M_j}e^{2\omega}f\,\mathrm dV_g=0\)，也就是
\(\int_{M_j}f\,\mathrm dV_{\widetilde g}=0\)。梯度能量在二维保持不变，但源项的总量仍须相对于新体积计算。
\end{solution}

% 自拟：正移位趋零的分支分解与误差估计，增加奇异极限的训练。
\Needspace{6\baselineskip}
\begin{exercise}
\label{b2:ex:26-vanishing-shift}
设 \((M,g)\) 为非空、紧致、无边界的 \(n\geq1\) 维光滑 Riemann 流形，连通分支为 \(M_1,\ldots,M_r\)。对实值 \(f\in L^2(M)\)，令
\[
 Pf=\sum_{j=1}^r\left(\frac1{\operatorname{Vol}_g(M_j)}
             \int_{M_j}f\,\mathrm dV_g\right)\mathbf1_{M_j}.
\]
对 \(\varepsilon>0\)，令 \(u_\varepsilon\in H^1(M)\) 为
\((-\Delta_g+\varepsilon)u_\varepsilon=f\) 的唯一弱解。证明
\(u_\varepsilon=\varepsilon^{-1}Pf+w_\varepsilon\)，其中 \(w_\varepsilon\) 在每个分支上平均值为零。设 \(u\) 是 \(-\Delta_gu=f-Pf\) 的分支零平均值解，证明存在仅依赖于 \((M,g)\) 的常数 \(C>0\)，使
\[
 \|w_\varepsilon-u\|_{H^1}
 \leq C\varepsilon\|f-Pf\|_2.
\]
由此判定 \(u_\varepsilon\) 在 \(H^1\) 中何时有有限极限。
\end{exercise}
\begin{solution}
在移位弱方程中取 \(v=\mathbf1_{M_j}\)，得到
\(\varepsilon\int_{M_j}u_\varepsilon=\int_{M_j}f\)，所以
\(w_\varepsilon=u_\varepsilon-\varepsilon^{-1}Pf\) 的每个分支平均值为零。常数部分的梯度消失，故
\[
 \int_Mg(\nabla_gw_\varepsilon,\nabla_gv)\,\mathrm dV_g
  +\varepsilon\int_Mw_\varepsilon v\,\mathrm dV_g
 =\int_M(f-Pf)v\,\mathrm dV_g
\]
对所有 \(v\in H^1(M)\) 成立。令 \(z=w_\varepsilon-u\)，减去 \(u\) 的方程，并取 \(v=z\)，得到
\[
 \|\nabla_gz\|_2^2+\varepsilon\|z\|_2^2
 =-\varepsilon\int_Muz\,\mathrm dV_g
 \leq\varepsilon C_P\|u\|_2\|\nabla_gz\|_2.
\]
因为 \(z\) 的分支平均值为零，这里使用的是分支 Poincaré 不等式。若梯度范数非零，约去它并丢掉左边第二项；为零时 \(z=0\)，同一估计仍成立。于是
\[
 \|z\|_{H^1}
 \leq\sqrt{1+C_P^2}\,\|\nabla_gz\|_2
 \leq\varepsilon C_P\sqrt{1+C_P^2}\,\|u\|_2.
\]
又由\cref{b2:cor:26-poisson-solvability}，有
\(\|u\|_2\leq C_P^2\|f-Pf\|_2\)，因此可取
\(C=C_P^3\sqrt{1+C_P^2}\)。这证明 \(w_\varepsilon\to u\) 于 \(H^1\)。若 \(Pf=0\)，则 \(u_\varepsilon=w_\varepsilon\to u\)；若 \(Pf\ne0\)，常数部分与 \(w_\varepsilon\) 在 \(L^2\) 中正交，故
\(\|u_\varepsilon\|_2^2\geq\varepsilon^{-2}\|Pf\|_2^2\to\infty\)，不可能有有限的 \(H^1\) 极限。
\end{solution}

% 自拟：由标量弱方程构造向量场的梯度投影，不调用微分形式 Hodge 理论。
\Needspace{6\baselineskip}
\begin{optionalexercise}
\label{b2:ex:26-gradient-projection}
设 \((M,g)\) 为非空、紧致、无边界的 \(n\geq1\) 维光滑 Riemann 流形，\(X\in L^2(TM)\) 为实切向量场，这里
\(\|X\|_2^2=\int_Mg(X,X)\,\mathrm dV_g<\infty\)。证明存在唯一的分支零平均值 \(u\in H^1(M)\)，使 \(Y=X-\nabla_gu\) 满足
\[
 \int_Mg(Y,\nabla_gv)\,\mathrm dV_g=0
 \quad(v\in H^1(M)).
\]
证明 \(X=\nabla_gu+Y\) 是 \(L^2(TM)\) 中的正交分解，并且 \(\nabla_gu\) 是 \(X\) 在全部 \(H^1\) 梯度中唯一的最佳 \(L^2\) 逼近。说明上式如何表达 \(\operatorname{div}_gY=0\) 的弱意义。
\end{optionalexercise}
\begin{solution}
定义
\(F(v)=\int_Mg(X,\nabla_gv)\,\mathrm dV_g\)。Cauchy--Schwarz 不等式给出
\(|F(v)|\leq\|X\|_2\|v\|_{H^1}\)，而各分支示性函数的梯度为零，所以 \(F(\mathbf1_{M_j})=0\)。由\cref{b2:cor:26-poisson-solvability}，存在唯一分支零平均值 \(u\)，使
\(\int g(\nabla_gu,\nabla_gv)=F(v)\) 对所有 \(v\in H^1\) 成立。移项即得 \(Y\) 的所需性质。

特别取 \(v=u\)，得到 \(\int g(Y,\nabla_gu)=0\)，所以
\(\|X\|_2^2=\|\nabla_gu\|_2^2+\|Y\|_2^2\)。对任意 \(w\in H^1(M)\)，还可取 \(v=u-w\)，从而
\[
 \|X-\nabla_gw\|_2^2
 =\|Y+\nabla_g(u-w)\|_2^2
 =\|Y\|_2^2+\|\nabla_g(u-w)\|_2^2.
\]
因此最小值为 \(\|Y\|_2\)，且取到最小值当且仅当 \(\nabla_gw=\nabla_gu\)。最佳梯度唯一，产生它的函数仍可相差分支常数。最后，弱散度按
\(\langle\operatorname{div}_gY,v\rangle=-\int g(Y,\nabla_gv)\)
定义为 \(H^{-1}(M)\) 元素，所列等式正是这个泛函等于零。
\end{solution}


% 自拟：平方单位分解下的能量恒等式，训练切断交叉项的精确抵消。
\Needspace{6\baselineskip}
\begin{exercise}[平方单位分解与能量局部化]\label{b2:ex:28-square-partition-energy}
设 \((M,g)\) 为非空紧致无边界光滑 Riemann 流形，\(\dim M=n\geq1\)。设 \(A:TM\to TM\) 是光滑、关于 \(g\) 对称且处处正定的实束自同态，\(c\in C^\infty(M;\mathbb R)\)、\(c\geq0\)，并令
\[
 \mathcal E(u)=\int_M\bigl(g(A\nabla_g u,\nabla_g u)+cu^2\bigr)\,\mathrm dV_g
 \qquad(u\in H^1(M;\mathbb R)).
\]
设有限个实光滑函数 \(\chi_1,\ldots,\chi_N\) 满足 \(\sum_{j=1}^N\chi_j^2=1\)。证明
\[
 \sum_{j=1}^N\mathcal E(\chi_j u)
 =\mathcal E(u)+\int_Mu^2\sum_{j=1}^N
       g(A\nabla_g\chi_j,\nabla_g\chi_j)\,\mathrm dV_g.
\]
由此证明，存在与 \(u\) 无关的 \(C\geq0\)，使
\[
 \mathcal E(u)\leq\sum_j\mathcal E(\chi_j u)
       \leq\mathcal E(u)+C\|u\|_{L^2(M)}^2.
\]
最后，从任意有限光滑单位分解 \(\eta_j\geq0\)、\(\sum_j\eta_j=1\) 构造具有相同支集控制的 \(\chi_j\)。
\end{exercise}
\begin{solution}
由 \(\sum_j\chi_j^2=1\) 求梯度，得到
\(\sum_j\chi_j\nabla_g\chi_j=0\)。将
\(\nabla_g(\chi_j u)=\chi_j\nabla_g u+u\nabla_g\chi_j\)
代入能量。含两个 \(\nabla_g u\) 的项求和后是原梯度能量；交叉项求和为
\[
 2u\,g\left(A\nabla_g u,\sum_j\chi_j\nabla_g\chi_j\right)=0;
\]
势项求和为 \(cu^2\sum_j\chi_j^2=cu^2\)。余下恰是题中列出的局部化误差。对 \(H^1\) 函数，所有乘积都可积，也可由光滑逼近通过到极限。

令 \(R=\sum_jg(A\nabla_g\chi_j,\nabla_g\chi_j)\)。正定性给出 \(R\geq0\)，光滑性与紧性给出 \(C=\max_MR<\infty\)，从而得到两个不等式。

对于给定的 \(\eta_j\)，令
\[
 s=\left(\sum_{j=1}^N\eta_j^2\right)^{1/2},\qquad
 \chi_j=\eta_j/s.
\]
由 Cauchy--Schwarz 不等式，\(1=(\sum_j\eta_j)^2\leq N\sum_j\eta_j^2\)，故 \(s\geq N^{-1/2}\)。因此 \(s\) 和 \(s^{-1}\) 都光滑，\(\sum_j\chi_j^2=1\)，且 \(\operatorname{supp}\chi_j=\operatorname{supp}\eta_j\)。这说明能量局部化可使用与坐标覆盖相容的平方单位分解。
\end{solution}

% 自拟：用统一椭圆估计证明解对系数参数的可微依赖，区别于正文的固定算子估计。
\Needspace{6\baselineskip}
\begin{optionalexercise}[解对系数参数的可微依赖]\label{b2:ex:28-coefficient-parameter}
设 \((M,g)\) 为非空紧致无边界光滑 Riemann 流形，\(\dim M=n\geq1\)，\(m\geq0\) 为整数。设光滑实束自同态 \(A,B:TM\to TM\) 均关于 \(g\) 对称，并且
\[
 g(AX,X)\geq\lambda|X|_g^2,
 \qquad |BX|_g\leq K|X|_g
 \quad(p\in M,\ X\in T_pM),
\]
其中 \(\lambda>0\)、\(K\geq0\)。选 \(\delta>0\) 满足 \(\delta K<\lambda/2\)，并取 \(c\in C^\infty(M;\mathbb R)\)、\(c\geq0\)、实值 \(f\in H^m(M)\)。对 \(|t|<\delta\)，令 \(u_t\in H^1(M;\mathbb R)\) 为唯一的弱解
\[
 L_tu_t=f,\qquad
 L_t=-\operatorname{div}_g((A+tB)\nabla_g)+(c+1).
\]
证明映射 \(t\mapsto u_t\) 从 \((-\delta,\delta)\) 到 \(H^{m+2}(M)\) 是连续可微的，并证明其导数 \(w_t\) 满足
\[
 L_tw_t=\operatorname{div}_g(B\nabla_g u_t).
\]
\end{optionalexercise}
\begin{solution}
对 \(|t|<\delta\)，主部的椭圆下界至少为 \(\lambda/2\)，势项 \(c+1\geq1\)。Lax--Milgram 方法给出每个 \(L_t\) 的唯一能量解。若 \(L_tv=h\)、\(h\in H^m(M)\)，以 \(v\) 测试便得
\[
 \|v\|_2^2\leq\int_Mhv\,\mathrm dV_g\leq\|h\|_2\|v\|_2,
 \qquad\|v\|_2\leq\|h\|_2.
\]
由\cref{b2:thm:28-global-higher}，\(v\in H^{m+2}(M)\)，且
\begin{equation}\label{b2:eq:28-ex-parameter-inverse}
 \|v\|_{H^{m+2}(M)}\leq C_m\|h\|_{H^m(M)}.
\end{equation}
常数可统一取于 \(|t|<\delta\)：有限坐标图保持不变，椭圆下界统一为 \(\lambda/2\)，而 \(A+tB\) 各个所需有限阶导数的上界由 \(A,B,\delta\) 控制。记 \(R_t=L_t^{-1}\)，则上一式说明 \(R_t:H^m\to H^{m+2}\) 一致有界。

置 \(P=-\operatorname{div}_g(B\nabla_g)\)。光滑系数的乘积法则说明 \(P:H^{m+2}\to H^m\) 有界，而且 \(L_t-L_s=(t-s)P\)。对任意 \(h\in H^m\)，将 \(L_t\) 作用于 \(R_th-R_sh\)，得到
\[
 R_t-R_s=-(t-s)R_tPR_s.
\]
各算子的定义域与值域正好依次匹配，所以该恒等式给出
\(\|R_t-R_s\|_{H^m\to H^{m+2}}\leq C|t-s|\)。特别地，\(u_t=R_tf\) 在 \(H^{m+2}\) 中连续。

再用上式取 \(s=t\)、将第一个参数改为 \(t+h\)，得
\[
 \frac{u_{t+h}-u_t}{h}=-R_{t+h}Pu_t
 \longrightarrow-R_tPu_t\quad\text{于 }H^{m+2}(M).
\]
因此 \(w_t=-R_tPu_t\)。已证的 \(R_t\) 算子范数连续性、\(u_t\) 连续性及 \(P\) 的有界性同时给出 \(w_t\) 连续。最后作用 \(L_t\)，有
\(L_tw_t=-Pu_t=\operatorname{div}_g(B\nabla_g u_t)\)，完成证明。
\end{solution}


% 第28章谱与模型增量习题；均为自拟，正文主线不调用此文件中的结论。
% 候选一：各向异性几何尺度、最小正特征值及重数。
\begin{exercise}\label{b2:ex:26-rectangular-torus}
设整数 \(n\geq2\)，\(a_1,\ldots,a_n>0\)。在 \(\mathbb R^n/(2\pi\mathbb Z)^n\) 上取平坦度量
\[
 g=\sum_{i=1}^n a_i^2\,\mathrm dx_i^2.
\]
求体积测度、\(-\Delta_g\) 的全部特征值及一组复 \(L^2\) 正交归一特征基。令 \(a_*=\max_i a_i\)，证明零均值函数的最佳 Poincaré 不等式为
\[
 \|u\|_{L^2(g)}\leq a_*\|\nabla_g u\|_{L^2(g)},
\]
并描述全部非零等号情形及最小正特征值的重数。最后在 \(n=2\)、\(a_1=2\)、\(a_2=1\) 时，求 \(-\Delta_gu=\cos(x_1+x_2)+\cos(2x_1)\) 的唯一零均值解。
\end{exercise}
\begin{solution}
度量矩阵为对角矩阵，逆矩阵的第 \(i\) 个对角元为 \(a_i^{-2}\)，体积密度为常数 \(\prod_i a_i\)。于是
\[
 \mathrm dV_g=\left(\prod_i a_i\right)\mathrm dx,\qquad
 -\Delta_g=-\sum_i a_i^{-2}\partial_{x_i}^2.
\]
令 \(V=(2\pi)^n\prod_i a_i\)。\cref{b2:prop:26-torus-spectrum}的 Fourier 完备性在测度乘以正常数后保持，故
\[
 E_k=V^{-1/2}e^{\mathrm ik\cdot x},\qquad
 \lambda_k=\sum_i\frac{k_i^2}{a_i^2}\quad(k\in\mathbb Z^n)
\]
给出所求基与全部特征值。不同 \(k\) 若有相同的这个数，就属于同一特征空间。

若 \(k\ne0\)，至少一个 \(|k_i|\geq1\)，所以 \(\lambda_k\geq a_*^{-2}\)。等号只能在 \(k=\pm e_i\) 且 \(a_i=a_*\) 时成立：任何额外的非零坐标都会增加正项，任何 \(|k_i|\geq2\) 也会严格增加该项。若达到最大长度的指标有 \(m\) 个，最小正特征值的实、复重数均为 \(2m\)。对零均值 \(u\)，谱能量公式给出
\[
 \|\nabla_g u\|_2^2=\sum_{k\ne0}\lambda_k|u_k|^2
 \geq a_*^{-2}\sum_{k\ne0}|u_k|^2=a_*^{-2}\|u\|_2^2.
\]
等号当且仅当所有非零系数都落在上述最小特征空间；实值时即
\[
 u(x)=\sum_{i:a_i=a_*}\bigl(\alpha_i\cos x_i+\beta_i\sin x_i\bigr),
\]
系数不全为零。它们也证明常数 \(a_*\) 最佳。

最后两个源项的特征值分别为 \(1/4+1=5/4\) 与 \(4/4=1\)，所以
\[
 u(x_1,x_2)=\frac45\cos(x_1+x_2)+\cos(2x_1).
\]
每一项平均值为零；直接代入核验方程，零均值唯一性来自 Poisson 可解性定理。
\end{solution}

% 候选二：二次多项式分层与非零频率共振。
\begin{exercise}\label{b2:ex:26-spherical-resonance}
设整数 \(n\geq1\)，\(B\) 为实对称 \((n+1)\times(n+1)\) 矩阵，\(b\in\mathbb R^{n+1}\)、\(c\in\mathbb R\)。在标准单位球面 \(S^n\) 上考虑
\[
 (-\Delta_{S^n}-n)u(\omega)=\omega^{\mathsf T}B\omega+b\cdot\omega+c.
\]
求存在 \(H^1(S^n;\mathbb R)\) 弱解的充要条件，并在条件满足时写出全部解。再给出满足 \(\int_{S^n}u\omega_i\,\mathrm d\sigma=0\)、\(1\leq i\leq n+1\) 的唯一解。说明仅要求 \(u\) 的平均值为零为何不能代替这组归一化。
\end{exercise}
\begin{solution}
令 \(\tau=\operatorname{tr}B/(n+1)\)，将源项分成
\[
 f=(c+\tau)+b\cdot\omega+Q(\omega),\qquad
 Q(\omega)=\omega^{\mathsf T}B\omega-\tau.
\]
常数、一次项、\(Q\) 分别属于零、一、二阶球谐空间，因为
\(x^{\mathsf T}Bx-\tau|x|^2\) 的欧氏 Laplace 为
\(2\operatorname{tr}B-2(n+1)\tau=0\)。三层的特征值分别为 \(0\)、\(n\)、\(2(n+1)\)。共振空间恰为一次球谐空间
\(\operatorname{span}\{\omega_1,\ldots,\omega_{n+1}\}\)。源项与这个空间正交的充要条件是 \(b=0\)：零、二阶与它正交，而一次项本身为零当且仅当 \(b=0\)。

条件满足时，全部解为
\[
 u(\omega)=-\frac{c+\tau}{n}+\frac{Q(\omega)}{n+2}+d\cdot\omega,
 \qquad d\in\mathbb R^{n+1}.
\]
这是逐层除以 \(-n\)、\(2(n+1)-n=n+2\)，再加上任意共振项的结果，所有项均光滑。由反射与坐标置换对称性，
\[
 \int_{S^n}\omega_i\omega_j\,\mathrm d\sigma
  =\frac{|S^n|}{n+1}\mathbf1_{\{i=j\}},
\]
故所给的 \(n+1\) 个归一化条件恰好迫使 \(d=0\)。所有一次球谐的平均值都为零，单独规定平均值无法确定 \(d\)；若 \(c+\tau\ne0\)，上述全部解还具有同一个非零平均值，此时强行要求零均值甚至会排除所有解。
\end{solution}

% 候选三：谱截断的误差控制与其最佳常数；训练数据正则性和解正则性的区别。
\begin{exercise}\label{b2:ex:26-spectral-truncation}
设 \((M,g)\) 为非空紧致无边界光滑 Riemann 流形，\(\dim M\geq1\)，有 \(r\) 个连通分支。取按重数递增排列的实正交归一特征基 \((e_j,\lambda_j)\)，其中 \(\lambda_1=\cdots=\lambda_r=0\)。对每个分支平均值为零的 \(f\in L^2(M;\mathbb R)\)，令 \(u\) 为零分支平均的 Poisson 解，并对整数 \(N\geq r\) 令
\[
 f_j=(f,e_j)_{L^2},\qquad u_N=\sum_{r<j\leq N}\frac{f_j}{\lambda_j}e_j.
\]
证明
\[
 \|u-u_N\|_2\leq\lambda_{N+1}^{-1}\|f\|_2,\qquad
 \|\nabla_g(u-u_N)\|_2\leq\lambda_{N+1}^{-1/2}\|f\|_2,
\]
且两个常数均最佳。再证明对每个固定 \(f\)，\(u_N\to u\) 于 \(H^2(M)\)，但不存在 \(\varepsilon_N\to0\)，使对所有此类 \(f\) 同时有
\[
 \|\Delta_g(u-u_N)\|_2\leq\varepsilon_N\|f\|_2.
\]
\end{exercise}
\begin{solution}
由谱展开与能量公式，
\[
 \|u-u_N\|_2^2=\sum_{j>N}\frac{|f_j|^2}{\lambda_j^2}
 \leq\lambda_{N+1}^{-2}\sum_{j>N}|f_j|^2,
\]
以及
\[
 \|\nabla_g(u-u_N)\|_2^2=\sum_{j>N}\frac{|f_j|^2}{\lambda_j}
 \leq\lambda_{N+1}^{-1}\sum_{j>N}|f_j|^2.
\]
尾和不大于 \(\|f\|_2^2\)，取平方根得到两式。取 \(f=e_{N+1}\) 时两式都为等号，因此常数最佳；即使 \(N\) 截在一个重特征空间中间，这个选择仍成立。

另一方面，
\[
 -\Delta_g(u-u_N)=\sum_{j>N}f_je_j,
\]
故其 \(L^2\) 范数对每个固定 \(f\) 趋于零。全局二阶估计给出
\[
 \|u-u_N\|_{H^2}
 \leq C\left(\left(\sum_{j>N}|f_j|^2\right)^{1/2}+\|u-u_N\|_2\right)
 \longrightarrow0.
\]
但若要求对所有单位球中的 \(f\) 一致成立，再取随 \(N\) 变化的 \(f=e_{N+1}\)，便有
\(\|\Delta_g(u-u_N)\|_2=\|f\|_2=1\)。于是任何统一的 \(\varepsilon_N\) 都至少为一。求解在零阶与一阶范数中压低高频误差，而把解再求两阶导数后，源项的高频尾部被完整恢复，只有对固定数据的收敛。
\end{solution}
